\section{Chapter~\ref{sec:gaba}. \nt{GLUT} and \nt{GABA} Together}

The chapter's logical and dynamical claims are derived in the chapter itself from linear analysis and Ohm's law; experiment shows that the circuits they rely on (the delayed veto, direct inhibition following excitation, stabilization by inhibition, the rhythm of the \nt{GLUT}--\nt{GABA} loop) really do operate in the brain.

{\footnotesize
\setlength{\tabcolsep}{3pt}\renewcommand{\arraystretch}{1.15}
\begin{longtable}{>{\raggedright}p{0.5cm}>{\raggedright}p{4.0cm}>{\raggedright}p{5.6cm}>{\raggedright}p{2.3cm}>{\raggedright\arraybackslash}p{1.7cm}}
\toprule
No. & Claim & Experiment and what was measured & Source & Status \\
\midrule\endfirsthead
\toprule
No. & Claim & Experiment and what was measured & Source & Status \\
\midrule\endhead
1 & A fifth to a quarter of cortical neurons are \nt{GABA} & GABA immunostaining in $10$ areas of monkey cortex: about $25\,\%$ of neurons in $8$ areas, less than $20\,\%$ in primary visual cortex. A review of the diversity of cortical inhibitory neurons & \cite{Hendry1987,Markram2004} & measured \\
2 & What inhibition does is largely decided by its landing site: dendrite, cell body, spike initiation site, or the terminal of another neuron's wire & Review: classes of cortical \nt{GABA} neurons differ in their target on the recipient. Simultaneous recording from the soma and dendrite of a pyramidal neuron: feedforward inhibition is much stronger at the soma than at the dendrites. Inhibition at the terminal of another neuron's wire in the spinal cord reduces transmitter release from one input line (review of measurements) & \cite{Markram2004,Pouille2001,Rudomin1999} & measured \\
3 & A sign change through an intermediary costs one delay, about a millisecond; inhibition arriving right after excitation narrows the coincidence window & Rat hippocampal pyramidal neurons in slices: inhibition through one intermediary follows direct excitation after a short delay, and the window in which two excitatory inputs sum to threshold is less than $2\ms$. On the dendrites, where this inhibition is weaker, the window is wider & \cite{Pouille2001} & measured \\
4 & The veto $a\wedge\bar b$~\eqref{eq:veto} with OR and a constant one is functionally complete (Proposition~\ref{prop:gabacomplete}) & Proof in the chapter. A classical result: a network of threshold elements with inhibitory inputs implements any logical expression. A claim about the model; no experiment & \cite{McCullochPitts1943} & theory \\
5 & A delayed veto gives a direction-of-motion detector with optimal speed $v^*=\Delta x/d$~\eqref{eq:vstar} & Rabbit retina: direction selectivity is created by inhibition that cancels excitation for motion in the ``null'' direction. Separate recording of the excitatory and inhibitory inputs of selective cells showed that starburst amacrine cells (\nt{GABA}) inhibit them asymmetrically, only with processes pointing to the ``null'' side. The formula for $v^*$ is the chapter's derivation & \cite{Barlow1965,Fried2002} & measured \\
6 & Shunting inhibition divides the voltage rather than subtracting~\eqref{eq:shunt} (Fig.~\ref{fig:shunt}) & Ohm's law for a divider of three conductances with the chloride equilibrium potential near rest; for a neuron that does not fire & derivation in the chapter & theory \\
7 & On the spike rate the shunt acts by subtraction, and division comes from inhibition together with balanced background noise & Model: in a firing neuron the current through the shunt is nearly constant, and the effect on rate is subtractive. Experiment: background excitatory and inhibitory conductances were injected into rat somatosensory cortex pyramidal neurons (slices) with a dynamic clamp; a simultaneous balanced increase of both divides the slope of the response with no noticeable shift in rate. Review: division by total activity occurs in many parts of the brain & \cite{HoltKoch1997,Chance2002,Carandini2012} & measured \\
8 & For $\beta>1$ mutual inhibition selects a single winner (Proposition~\ref{prop:wta}, \eqref{eq:wta}) & The eigenvalues $-(1\pm\beta)/\tau$ are derived in the chapter. The rates $0.73$ and $0.53$ at $\beta=0.5$ are the fixed point $r_{1,2}=(I_{1,2}-\beta I_{2,1})/(1-\beta^2)$; there is no script in \texttt{models/}. No direct experiment is cited in the chapter & derivation in the chapter & theory \\
9 & Memory reset by a signal through \nt{GABA}; two groups with mutual inhibition form a latch & Follows from Fig.~\ref{fig:bistable} and Proposition~\ref{prop:wta}. No experiment & derivation in the chapter & theory \\
10 & The equilibrium of the \nt{GLUT}--\nt{GABA} loop is stable if the inhibition is strong enough and fast enough (Proposition~\ref{prop:ei}) & Two-population model~\eqref{eq:ei}; conditions (i) and (ii) are the determinant and trace of its matrix, derived in the chapter & \cite{WilsonCowan1972} & theory \\
11 & For $w_{EE}=1.5$ and $w_{EI}w_{IE}=2$, fast inhibition ($\tau_I=2\ms$) holds $E^*=0.67h$, and slow inhibition ($30\ms$) drives oscillations (Fig.~\ref{fig:ei}) & Numerical solution of~\eqref{eq:ei}, script \texttt{figures/make\_ei.py} & calculation & model \\
12 & The cortex works in an inhibition-stabilized regime; the signature: push the \nt{GABA} neurons, and their rate drops~\eqref{eq:isn} & Theory predicted the paradoxical response. Experiment: intracellular recordings in cat primary visual cortex show that when the response is suppressed by a surrounding stimulus, the inhibitory conductance does not rise but falls together with the excitatory one. Optogenetic excitation of PV neurons in mouse visual, somatosensory, and motor cortex lowers the rate of inhibitory neurons, including without a stimulus and under light anesthesia. The coefficient $-1/3$ for the numbers of Fig.~\ref{fig:ei} is the chapter's derivation & \cite{Tsodyks1997isn,Ozeki2009,Sanzeni2020} & measured \\
13 & The loop ``\nt{GLUT} excites \nt{GABA}, \nt{GABA} inhibits \nt{GLUT}'' generates a fast rhythm with a period of $12$--$50\ms$ & Optogenetic stimulation of fast \nt{GABA} neurons in mouse somatosensory cortex at $8$--$200$\,Hz selectively enhances the gamma rhythm ($20$--$80$\,Hz, period $12$--$50\ms$); stimulation of \nt{GLUT} neurons enhances only slower rhythms & \cite{Cardin2009,Buzsaki2012} & measured \\
\bottomrule
\end{longtable}}
